\par\SourceLeaf{418}%
\begin{center}3.\end{center}
\noindent\hspace*{1.2em}Wenn wir die Lage der Punkte in der Ebene auf zwei auf einander senkrechte Coordinatenaxen beziehen, denen die Richtungen $0$ und $90^{0}$ entsprechen, so dass die erste Coordinate den nach der Richtung der ersten Axe gemessenen Abstand von der zweiten Axe, die zweite Coordinate hingegen den nach der Richtung der zweiten Axe gemessenen Abstand von der ersten Axe vorstellen; wenn ferner $x,y$ unbestimmt die Coordinaten eines Punktes der krummen Linie, $s$ deren Länge von einem beliebigen Anfangspunkte an bis zu diesem gezählt, $\varphi$ die Richtung der Tangente an demselben, $r$ den Krümmungshalbmesser ausdrücken, so wird
\[
\begin{aligned}
dx&=\cos\varphi.ds\\
dy&=\sin\varphi.ds\\
r&=\frac{ds}{d\varphi}.
\end{aligned}
\]

Wird die Natur der krummen Linie durch die Gleichung $V=0$ vorgestellt, wo $V$ eine Function von $x$ und $y$ ist, und setzt man
\[
dV=pdx+qdy,
\]
so wird in der krummen Linie
\[
pdx+qdy=0
\]
sein; also
\[
p\cos\varphi+q\sin\varphi=0,
\]
woraus also
\[
\operatorname{tang}\varphi=-\frac pq.
\]
Es ist ferner
\[
\cos\varphi.dp+\sin\varphi.dq-(p\sin\varphi-q\cos\varphi)d\varphi=0.
\]
Setzt man also nach einem bekannten Lehrsatz
\[
\begin{aligned}
dp&=Pdx+Qdy\\
dq&=Qdx+Rdy,
\end{aligned}
\]
